\documentclass[11pt,a4paper]{article}

% Packages for styling and functionality
\usepackage[utf8]{inputenc}
\usepackage[T1]{fontenc}
\usepackage[english]{babel}
\usepackage{amsmath, amssymb, amsfonts}
\usepackage{graphicx}
\usepackage{hyperref}
\usepackage{geometry}
\usepackage{fancyhdr}
\usepackage{titlesec}
\usepackage{booktabs}
\usepackage{xcolor}

% Page geometry
\geometry{margin=2.5cm}

% Header and footer
\pagestyle{fancy}
\fancyhf{}
\rhead{LLM VRAM Calculator}
\lhead{Technical Documentation}
\rfoot{Page \thepage}

% Title styling
\titleformat{\section}{\large\bfseries}{\thesection}{1em}{}[\titlerule]

\title{
    \vspace{-2cm}
    \textbf{LLM VRAM Calculator} \\
    \large Technical Documentation and Methodology
}
\author{Croix-Rouge luxembourgeoise}
\date{\today}

\begin{document}

\maketitle
\tableofcontents
\newpage

\section{Introduction}
The \textbf{LLM VRAM Calculator} is a professional estimation tool designed to help researchers and engineers accurately project the hardware requirements for deploying Large Language Models (LLMs). 

As models grow in size and complexity (e.g., Mixture-of-Experts, GQA/MQA architectures), understanding memory constraints and performance bottlenecks is critical for cost-effective deployment.

\section{Core Features}
\begin{itemize}
    \item \textbf{Hugging Face Integration}: Direct metadata import from the HF Hub.
    \item \textbf{Architecture Awareness}: Support for Dense and MoE models.
    \item \textbf{Quantization Modeling}: Support for GGUF, GPTQ, AWQ, EXL2, and NF4.
    \item \textbf{Performance Estimation}: Roofline modeling for latency and throughput.
    \item \textbf{Cost Analysis}: Power consumption and CO\textsubscript{2} emission estimates.
\end{itemize}

\section{Methodology}
\subsection{Memory Estimation}
The total VRAM consumption is calculated as the sum of weights, KV cache, and activation memory:
\begin{equation}
    VRAM_{total} = VRAM_{weights} + VRAM_{KV\_cache} + VRAM_{activations}
\end{equation}

\subsubsection{Weights Memory}
For a model with $P$ parameters and a quantization bit-depth $b$:
\begin{equation}
    VRAM_{weights} = P \times \frac{b}{8} \times 1.05 \text{ (overhead buffer)}
\end{equation}

\subsubsection{KV Cache Memory}
For Grouped-Query Attention (GQA):
\begin{equation}
    VRAM_{KV} = 2 \times L \times H_{kv} \times D \times S \times \frac{b_{kv}}{8}
\end{equation}
Where:
\begin{itemize}
    \item $L$: Number of layers
    \item $H_{kv}$: Number of KV heads
    \item $D$: Head dimension
    \item $S$: Sequence length (Context)
    \item $b_{kv}$: Precision bits for KV cache
\end{itemize}

\section{User Guide}
\subsection{Importing from Hugging Face}
To import a model, use the search bar to find the repository ID (e.g., \texttt{meta-llama/Llama-3.1-8B}). The tool will automatically fetch the configuration files (\texttt{config.json}, \texttt{generation\_config.json}) to populate the calculator.

\subsection{Configuring Quantization}
Select the desired quantization format from the dropdown. Note that different formats (e.g., GGUF vs. EXL2) have different overheads and hardware compatibility.

\section{Conclusion}
This tool aims to bridge the gap between theoretical model size and practical deployment constraints, enabling informed decisions on hardware procurement and optimization strategies.

\end{document}
